\documentclass[11pt]{article}
% ======================================================================
%  examstyle.tex — EPFL-style exam layout for the weekly Time Series exams
%  Shared by every Exam_Week_NN(.tex / _Solutions.tex). Compiles with pdflatex.
% ======================================================================
\usepackage[utf8]{inputenc}
\usepackage[T1]{fontenc}
\usepackage[english]{babel}
\usepackage[a4paper,margin=2.4cm]{geometry}
\usepackage{amsmath,amssymb,amsthm,mathtools}
\usepackage{bm}
\usepackage{enumitem}
\usepackage{array}

\setlength{\parindent}{0pt}
\setlength{\parskip}{0.5em}
\emergencystretch=3em
\allowdisplaybreaks[1]

% ---------- math macros (same names as the rest of the project) ----------
\newcommand{\E}{\mathbb{E}}
\DeclareMathOperator{\Var}{Var}
\DeclareMathOperator{\Cov}{Cov}
\DeclareMathOperator{\Corr}{Corr}
\DeclareMathOperator*{\argmin}{arg\,min}
\DeclareMathOperator{\MSE}{MSE}
\DeclareMathOperator{\trace}{tr}
\newcommand{\Reals}{\mathbb{R}}
\newcommand{\Ints}{\mathbb{Z}}
\newcommand{\Comp}{\mathbb{C}}
\newcommand{\Nat}{\mathbb{N}}
\newcommand{\Prob}{\mathbb{P}}
\newcommand{\Tr}{^{\mathsf{T}}}
\newcommand{\conj}[1]{\overline{#1}}
\newcommand{\abs}[1]{\left\lvert #1 \right\rvert}
\newcommand{\norm}[1]{\left\lVert #1 \right\rVert}
\newcommand{\ind}{\mathbf{1}}
\newcommand{\eps}{\varepsilon}
\newcommand{\diff}{\nabla}
\newcommand{\acvs}{\gamma}
\newcommand{\acf}{\rho}
\newcommand{\pacf}{\alpha}
\newcommand{\sdf}{\mathcal{S}}
\newcommand{\Bsh}{B}
\newcommand{\iid}{\overset{\text{iid}}{\sim}}

\setlist[enumerate]{leftmargin=2em,topsep=2pt,itemsep=4pt}
\setlist[itemize]{leftmargin=1.6em,topsep=2pt,itemsep=2pt}

\newcounter{exo}

% ---------- exam cover header :  \examheader{exam title}{duration} ----------
\newcommand{\examheader}[2]{%
  \thispagestyle{empty}
  \begin{center}
    {\Large\bfseries Time Series}\\[3pt]
    Spring Semester 2026, EPFL\\
    \'Ecole Polytechnique F\'ed\'erale de Lausanne\\
    Prof.\ Dr.\ Sofia C.\ Olhede\\[7pt]
    \hrule height 0.8pt
    \vspace{9pt}
    {\large\bfseries Practice Exam --- #1}\\[4pt]
    {\normalsize Duration: #2 \quad$\cdot$\quad No calculator \quad$\cdot$\quad Answer on paper}\\
    \vspace{8pt}
    \hrule height 0.8pt
  \end{center}
  \vspace{14pt}
  \begin{tabular}{@{}p{8.2cm}@{}p{8cm}@{}}
    Name:\enspace\hrulefill & Surname:\enspace\hrulefill
  \end{tabular}\\[14pt]
  SCIPER (Student Card Number):\enspace\hrulefill
  \vspace{16pt}
}

% ---------- points table (fixed: 4 problems) ----------
\newcommand{\pointstable}{%
  \begin{center}
  \renewcommand{\arraystretch}{1.5}
  \begin{tabular}{|p{5cm}|p{4cm}|}
    \hline
    \textbf{Exercise} & \textbf{Points}\\\hline
    1 & \\\hline
    2 & \\\hline
    3 & \\\hline
    4 & \\\hline
    \textbf{Total} & \\\hline
  \end{tabular}
  \end{center}
}

% ---------- problem heading (exam: one problem per page) ----------
\newcommand{\problem}[1]{%
  \clearpage
  \stepcounter{exo}%
  \begin{center}\textbf{\large Problem \theexo\ (#1 points)}\end{center}
  \vspace{4pt}
}
% ---------- answer space: fill the statement page + one blank answer page ----------
% Put \answerpage at the END of each problem so every exercise gets its own page(s).
\newcommand{\answerpage}{\par\vfill\newpage\null\newpage}

% ---------- solutions document ----------
\newcommand{\solheader}[1]{%
  \begin{center}
    {\Large\bfseries Time Series --- Solutions}\\[3pt]
    {\large Practice Exam --- #1}\\[2pt]
    EPFL, MATH-342 \quad$\cdot$\quad Prof.\ S.\ C.\ Olhede\\[5pt]
    \hrule height 0.8pt
  \end{center}
  \vspace{10pt}
}
\newcommand{\solproblem}[1]{%
  \bigskip
  \stepcounter{exo}%
  {\large\bfseries Problem \theexo\ (#1 points)}\par\nobreak\vspace{2pt}
}
% Rappel de l'énoncé avant la solution (dans les corrigés)
\newcommand{\statementstart}{\par\vspace{1pt}{\bfseries\itshape Statement.}\par\nobreak\begingroup\small\itshape}
\newcommand{\statementend}{\par\endgroup\vspace{5pt}\hrule\vspace{6pt}{\bfseries Solution.}\par\nobreak\vspace{2pt}}

\begin{document}

\solheader{Week 1: Stationarity, Autocovariance \& Ergodicity}

% ---------------------------------------------------------------
\solproblem{5}
\statementstart
\textbf{(a)} State the definition of \emph{weak (second-order) stationarity} and of \emph{strict stationarity} for a process $\{X_t\}_{t\in\Ints}$. \hfill(2 pts)

\textbf{(b)} Let $Y_t\iid\mathcal N(0,\sigma_Y^2)$ and let $a,b,c$ be constants. For each process, decide whether it is weakly and/or strictly stationary, and give its mean and ACVS:
\begin{enumerate}[label=(\roman*)]
  \item $X_t = a + bY_t + cY_{t-1}$;
  \item $X_t = Y_0\cos(ct)$.
\end{enumerate}\hfill(2 pts)

\textbf{(c)} Give a process that is strictly stationary but \emph{not} weakly stationary, and explain why. \hfill(1 pt)
\statementend

\textbf{(a)} $\{X_t\}$ is \emph{weakly (second-order) stationary} if all first- and second-order moments are finite and shift-invariant: $\E[X_t]=\mu$ (constant), $\Var(X_t)=\sigma^2<\infty$ (constant), and $\Cov(X_t,X_{t+\tau})$ depends on the lag $\tau$ only (not on $t$). It is \emph{strictly stationary} if for every $n$ and all $t_1,\dots,t_n$, the joint law of $(X_{t_1},\dots,X_{t_n})$ equals that of $(X_{t_1+\tau},\dots,X_{t_n+\tau})$ for all $\tau$.

\textbf{(b)(i)} $\E[X_t]=a$. By independence of the $Y_t$,
\[
\acvs_\tau=\begin{cases}(b^2+c^2)\sigma_Y^2,&\tau=0,\\ bc\,\sigma_Y^2,&\tau=\pm1,\\ 0,&\abs\tau\ge2.\end{cases}
\]
Depends on $\tau$ only and is finite $\Rightarrow$ weakly stationary; a linear combination of jointly Gaussian variables is Gaussian $\Rightarrow$ also strictly stationary.

\textbf{(b)(ii)} $\E[X_t]=0$ but $\Cov(X_t,X_{t+\tau})=\sigma_Y^2\cos(ct)\cos\!\big(c(t+\tau)\big)$ \emph{depends on $t$} $\Rightarrow$ \emph{not} stationary (neither weakly nor strictly).

\textbf{(c)} An i.i.d.\ sequence of Cauchy (Student-$t_1$) variables: i.i.d.\ hence strictly stationary, but it has no finite mean/variance, so the finite-variance requirement of weak stationarity fails. Thus strict $\not\Rightarrow$ weak.

% ---------------------------------------------------------------
\solproblem{5}
\statementstart
Let $\{X_t\}$ be a second-order stationary process with ACVS $\acvs_\tau$, \emph{independent} of the white noise $\{\eps_t\}$.

\textbf{(a)} Determine the ACVS of $Z_t = X_t + \eps_t$. \hfill(2 pts)

\textbf{(b)} Let $W_t = \theta_1\eps_t + \theta_2\eps_{t-1}$. Compute its ACVS $\acvs^W_\tau$ for all $\tau$ and its autocorrelation $\acf_1$, and prove that $\abs{\acf_1}\le\tfrac12$. \hfill(3 pts)
\statementend

\textbf{(a)} Independence kills all cross-covariances $\Cov(X_s,\eps_u)=0$, and the white noise contributes variance only at lag $0$:
\[
\acvs^Z_\tau=\acvs_\tau+\sigma^2\,\ind_{\{\tau=0\}}.
\]

\textbf{(b)} By bilinearity and $\Cov(\eps_s,\eps_u)=\sigma^2\ind_{\{s=u\}}$ (only equal noise indices survive),
\[
\acvs^W_\tau=\begin{cases}(\theta_1^2+\theta_2^2)\sigma^2,&\tau=0,\\ \theta_1\theta_2\,\sigma^2,&\tau=\pm1,\\ 0,&\abs\tau\ge2,\end{cases}
\qquad
\acf_1=\frac{\theta_1\theta_2}{\theta_1^2+\theta_2^2}.
\]
By AM--GM, $\abs{\theta_1\theta_2}\le\tfrac12(\theta_1^2+\theta_2^2)$, hence $\abs{\acf_1}\le\tfrac12$, with equality iff $\abs{\theta_1}=\abs{\theta_2}$. (No process of this MA(1)-type can exceed $\abs{\acf_1}=\tfrac12$.)

% ---------------------------------------------------------------
\solproblem{5}
\statementstart
\textbf{(a)} Prove that the autocovariance sequence $\{\acvs_\tau\}$ of a stationary process is \emph{positive semi-definite}: for every $n\ge1$ and all $a_1,\dots,a_n\in\Reals$,
\[
\sum_{j=1}^{n}\sum_{k=1}^{n} a_j a_k\,\acvs_{j-k}\ \ge\ 0 .
\]\hfill(2 pts)

\textbf{(b)} Deduce the symmetry $\acvs_{-\tau}=\acvs_\tau$ and the bound $\abs{\acvs_\tau}\le\acvs_0$. \hfill(1 pt)

\textbf{(c)} Is the sequence $\acvs_0 = 1,\ \acvs_{\pm1} = 0.8,\ \acvs_\tau = 0\ (\abs\tau\ge2)$ a valid ACVS of some stationary process? Justify. \hfill(2 pts)
\statementend

\textbf{(a)} Let $W=\sum_{j=1}^n a_j X_j$. A variance is non-negative, so
\[
0\le\Var(W)=\sum_{j=1}^n\sum_{k=1}^n a_j a_k\Cov(X_j,X_k)=\sum_{j=1}^n\sum_{k=1}^n a_j a_k\,\acvs_{j-k}.
\]

\textbf{(b)} Symmetry: $\acvs_{-\tau}=\Cov(X_t,X_{t-\tau})=\Cov(X_{t-\tau},X_t)=\acvs_\tau$ (real case). Bound: by Cauchy--Schwarz $\abs{\acvs_\tau}=\abs{\Cov(X_t,X_{t+\tau})}\le\sqrt{\Var(X_t)\Var(X_{t+\tau})}=\acvs_0$.

\textbf{(c)} \textbf{No.} A valid ACVS must be positive semi-definite, equivalently have a non-negative spectral density. Here
\[
\sdf(f)=\sum_\tau\acvs_\tau e^{-2\pi i f\tau}=1+1.6\cos(2\pi f),
\]
which equals $1-1.6=-0.6<0$ at $f=\tfrac12$. The spectral density goes negative $\Rightarrow$ the sequence is not p.s.d.\ $\Rightarrow$ not a valid ACVS. (Already $\abs{\acf_1}=0.8>\tfrac12$ is impossible for an MA(1).)

% ---------------------------------------------------------------
\solproblem{5}
\statementstart
For observations $X_1,\dots,X_n$ recall the two sample-autocovariance estimators (for $\abs\tau\le n-1$, and $0$ otherwise):
\[
\tilde\acvs_\tau=\frac{1}{n-\abs\tau}\sum_{t=1}^{n-\abs\tau}(X_t-\bar X)(X_{t+\abs\tau}-\bar X),\qquad
\hat\acvs_\tau=\frac{1}{n}\sum_{t=1}^{n-\abs\tau}(X_t-\bar X)(X_{t+\abs\tau}-\bar X).
\]

\textbf{(a)} Which estimator is biased? Which one is positive semi-definite? \hfill(1 pt)

\textbf{(b)} For Gaussian white noise (known mean) one has, for $\tau\ne0$, $\Var(\hat\acvs_\tau)=\tfrac{n-\abs\tau}{n^2}\sigma^4$ and $\Var(\tilde\acvs_\tau)=\tfrac{1}{n-\abs\tau}\sigma^4$. Which estimator has the smaller mean-squared error, and why? \hfill(2 pts)

\textbf{(c)} Show that if $\sum_{\tau}\abs{\acvs_\tau}<\infty$ then $\Var(\bar X)\to0$ as $n\to\infty$, so $\bar X$ is a consistent (mean-ergodic) estimator of $\mu=\E[X_t]$. \hfill(2 pts)
\statementend

\textbf{(a)} $\tilde\acvs_\tau$ is \emph{unbiased} (known mean) but not always positive semi-definite; $\hat\acvs_\tau$ is \emph{biased} ($\E[\hat\acvs_\tau]=\tfrac{n-\abs\tau}{n}\acvs_\tau$) but positive semi-definite --- which is why it is the default.

\textbf{(b)} For white noise both are unbiased, so $\MSE=\Var$. Now
\[
\Var(\hat\acvs_\tau)=\frac{n-\abs\tau}{n^2}\sigma^4=\Big(\tfrac{n-\abs\tau}{n}\Big)^2\frac{\sigma^4}{n-\abs\tau}=\Big(\tfrac{n-\abs\tau}{n}\Big)^2\Var(\tilde\acvs_\tau)\le\Var(\tilde\acvs_\tau).
\]
So the \emph{biased} estimator $\hat\acvs_\tau$ has the smaller MSE. (Lesson: the goal is low MSE and positive semi-definiteness, not unbiasedness.)

\textbf{(c)} $\E[\bar X]=\mu$, and grouping by lag,
\[
\Var(\bar X)=\frac1{n^2}\sum_{i,j=1}^n\acvs_{j-i}=\frac1n\sum_{\tau=-(n-1)}^{n-1}\Big(1-\tfrac{\abs\tau}{n}\Big)\acvs_\tau .
\]
By Ces\`aro summability, $\sum_\tau\big(1-\tfrac{\abs\tau}{n}\big)\acvs_\tau\to\sum_{\tau=-\infty}^\infty\acvs_\tau=:C(\acvs)$, finite since $\sum_\tau\abs{\acvs_\tau}<\infty$. Hence $n\Var(\bar X)\to C(\acvs)$, i.e. $\Var(\bar X)=O(1/n)\to0$, so $\bar X\to\mu$ in mean square, therefore in probability: $\bar X$ is consistent (the process is mean-ergodic).

\end{document}
